\section{Introduction}

A type-theoretic library can grow in two very different ways.  In the first,
a user defines a new transparent object inside an already fixed foundation.  A
new datatype of trees, for example, need not explain how it interacts with every
old construction in the library before it can be used.  Its dependencies are
local, and the surrounding theory can mostly ignore it.

The second kind of growth is less forgiving.  A foundation can be extended by a
new sealed structural layer: a modality, a universe former, a promoted interface
package, or some other globally acting operation whose public behavior becomes
part of the exported core.  Such an extension is not merely another transparent
term.  Once sealed, it advertises a stable public interface.  It must therefore
explain how it acts on the active structure already present: old type formers,
transport operations, eliminators, computation rules, paths, equivalences, and
comparison data.  This explanatory residue is what we call \emph{public trace}.

The question of this paper is:
\begin{quote}
When a cubical type-theoretic foundation is extended by a sealed, globally acting
structural layer, how much public integration trace must remain primitive after
normalization and opacity?
\end{quote}

The answer proved here is a depth-two theorem for one fixed formal calculus.  In
that calculus, unary trace is not enough: univalence supplies genuine binary
naturality obligations.  But trace above binary arity is not primitive: the
higher structural obligations generated by sealing are cubical open-box extension
problems, and cubical composition and filling compute their missing faces and
fillers.  Thus higher exact obligation factors remain visible in the semantics,
but they disappear from minimal opaque public signatures.

A useful slogan is:
\begin{quote}
Binary trace is necessary; higher structural trace is cubically payable.
\end{quote}

This slogan is deliberately about \emph{structural integration trace}, not about
all higher homotopy in a cubical theory.  The theorem does not say that higher
paths vanish, that user-supplied higher payload is forbidden, or that ordinary
libraries grow according to a fixed recurrence.  It says that, for the sealed
extension calculus studied below, the additional public trace forced by sealing a
globally acting structural layer stabilizes at binary trace.

\subsection{Payload, trace, and the first obstruction}

The first distinction is between \emph{payload} and \emph{trace}.  The payload is
what the new sealed layer contributes: a new operation, modality, universe rule,
or promoted structural package.  The trace is the public evidence that this
payload interacts correctly with the already exported interface.

Consider a small old interface containing two objects \(A\) and \(B\), together
with a comparison
\[
  p : A \simeq B
\]
or, in a univalent presentation, a path obtained from such an equivalence.  Now
seal a new globally acting operation \(X\).  There are several levels of public
obligation.

First, there is unary trace: the extension must say how \(X\) acts on individual
old generators such as \(A\) and \(B\).  This is objectwise information.  Second,
there is binary trace: the extension must say that the action of \(X\) respects
old comparison data such as \(p\).  This is not reducible to two unrelated unary
actions, because \(p\) is part of the old public interface.  A sealed action that
acts correctly on \(A\) and \(B\) separately may still fail to commute with the
public equivalence between them.

The basic lower-bound example is the two-point type.  Univalence turns the
nontrivial boolean swap equivalence into a public path.  A candidate sealed
action may be well formed objectwise but fail the naturality square along this
swap path.  The identity action passes the check; an action that ignores the
swap can fail it.  Hence unary trace alone cannot decide admissibility in
general.  Binary trace is a real public obligation, not an accounting artifact.

This is the first half of the depth-two story:
\[
  \text{depth one is insufficient.}
\]
The second half explains why the story stops at two.

\subsection{Why higher structural obligations are different}

At first sight, sealing seems to create an infinite tower.  If unary action must
respect binary comparison data, should binary comparison data itself not satisfy
ternary compatibility laws?  Should those laws not satisfy quaternary laws, and
so on?  This is the familiar shape of coherence problems.

Cubical foundations change the accounting.  Their primitive operations are built
to solve open-box extension problems.  Given a partial boundary of a cube and a
compatible base, cubical composition and filling supply the missing face and the
interior filler.  In a CCHM-style core this is the role of the Kan operations,
together with transport and the interval structure.  The point of the present
paper is that the higher structural obligations generated by sealing have exactly
this open-box form.

The relevant higher obligation is not an arbitrary closed higher-dimensional
choice.  It is a \emph{marginal extension problem}.  The lower-dimensional faces
come from the already exported unary and binary trace, together with normalized
public signatures of previous sealed layers.  What appears as a ternary or higher
coherence field is therefore represented as an open box: a partial boundary, a
compatibility condition, a missing remote face, and a filler.  In the fixed
calculus, this open-box extension object is contractible and substitution-stable.
Its canonical point is computed using cubical composition and filling.

This is the key semantic mechanism.  Higher structural obligations are not
ignored.  They remain present in the exact realized obligation object.  But they
are no longer primitive public data, because the calculus computes them from the
lower-arity boundary.  Equivalently, an explicit higher structural field in a
surface presentation can be replaced by a cubical term built from the displayed
boundary, and then erased from the minimal public signature.

Thus the main theorem has two levels.  At the level of exact realized obligation
objects, higher factors are present but contractible.  At the level of normalized
opaque public signatures, higher structural fields are absent from the primitive
trace count.

\subsection{The fixed calculus proved here}

Let \(C_{\mathrm{ext}}\) be the sealed-extension calculus defined in this paper.
It consists of a CCHM-style cubical core with univalence and Kan composition,
globally acting extension declarations, opaque public signatures, and a raw
adequacy bridge connecting surface presentations to normalized trace objects.
The definitions of \(C_{\mathrm{ext}}\) are part of the theorem object.  The paper
does not claim that this is the only possible formalization of sealed cubical
extension, nor that every cubical proof assistant program automatically
interprets into it.

For an admissible candidate \(X\), write \(\mathrm{Real}_k(X)\) for the exact
realized structural obligation object obtained when public trace is allowed up
to arity \(k\).  Write \(\mu\) for the primitive public trace cost of a normalized
opaque signature.  The central result is:
\begin{quote}
\textbf{Main theorem for \(C_{\mathrm{ext}}\).}
For every admissible candidate \(X\), the exact realized obligation objects
stabilize at binary trace:
\[
  \mathrm{Real}_k(X) \simeq \mathrm{Real}_2(X)
  \qquad (k \geq 2).
\]
Moreover, every \(\mu\)-minimal normalized public signature contains no primitive
structural trace field of arity greater than two.  Unary trace is insufficient in
general, as witnessed by the two-point swap obstruction.
\end{quote}

The theorem should be read as an exact-depth statement for a fixed calculus:
\[
  \text{binary trace is necessary, and binary trace is sufficient for primitive
  public structural integration.}
\]
The sufficiency direction is proved by the structural horn-to-open-box theorem,
the open-box contractibility theorem, and the computational replacement theorem
for explicit higher structural fields.  The necessity direction is proved by the
univalent boolean-swap obstruction.

This separation is important.  The paper does not prove contractibility of
arbitrary higher Homotopy Type Theory filler spaces.  It proves contractibility
of the theorem-facing structural open-box extension fibers generated by the
sealed-extension calculus.  Nor does it prove that every external cubical surface
language is adequate for \(C_{\mathrm{ext}}\).  Instead it isolates an adequacy
criterion: a surface calculus or cubical foundation inherits the theorem when it
admits a faithful interpretation preserving semantic admissibility,
primitive-versus-derived trace status, historical support, horn-computational
replacement, and public trace cardinality.

\subsection{From arity depth to chronological support}

Depth as arity is not the same thing as chronological locality.  A binary
comparison could, in principle, compare a new layer directly with a very old
layer.  Such a comparison still has arity two, but it has unbounded historical
reach.  Therefore the depth-two theorem alone does not imply a two-step
recurrence.

The calculus contains a separate recent-history theorem.  Under
factorization-complete public export, surviving primitive binary comparisons with
remote history factor through the normalized public signatures of the two most
recent layers.  In that setting the arity-depth theorem becomes a two-layer
chronological window.  Higher remote interactions are not counted as primitive;
they are mediated by the already exported public signatures.

Only after this semantic factorization does the counting theorem enter.  If a
sealed sequence is fully coupled to the active primitive basis, and if the
bundled per-site convention is used, then the primitive trace cost satisfies
\[
  \mu_{n+1} = \mu_n + \mu_{n-1} + \kappa_n + \kappa_{n-1},
\]
where \(\kappa_n\) is the payload contribution at layer \(n\).  Sparse footprints
satisfy the corresponding sparse formula and remain sparse.  The Fibonacci
statement is only the constant-payload endpoint: when \(\kappa_n=c\), the shifted
sequence \(U_n=\mu_n+2c\) follows the Fibonacci recurrence.

This order of explanation matters.  The recurrence is not the main phenomenon.
It is the arithmetic shadow of three prior facts: higher structural trace is
cubically derived, binary trace has a recent-history factorization, and the
chosen extension regime is fully coupled to the active basis.  Ordinary
transparent library growth need not satisfy those hypotheses and should not be
expected to display Fibonacci behavior.

\subsection{What the paper proves, and what remains conditional}

The contribution of the paper is a theorem stack for the fixed calculus
\(C_{\mathrm{ext}}\), together with a transfer interface for broader systems.

First, the paper defines a raw sealed-extension calculus that separates payload
from public trace, primitive fields from derived horn-computational fields, exact
realized obligation objects from \(\mu\)-minimal public signatures, and arity
depth from chronological support.

Second, it proves that every higher structural obligation generated by sealing is
represented as a theorem-facing open-box extension object.  These open-box
extension fibers are contractible and substitution-stable in the fixed cubical
core.  This is the semantic reason that higher exact factors do not contribute
primitive public trace.

Third, it proves a computational replacement theorem.  An explicit higher
structural field in a presentation elaborates to a cubical term, built from
composition, filling, and transport, using only its lower-arity boundary data.
After normalization, the explicit field is absent from the \(\mu\)-minimal public
signature.

Fourth, it proves exact structural integration depth two for \(C_{\mathrm{ext}}\):
\(\mathrm{Real}_k(X) \simeq \mathrm{Real}_2(X)\) for \(k\geq 2\), with no
primitive public trace above binary arity, while the swap obstruction shows that
unary trace alone is insufficient.

Fifth, it proves the downstream counting results: sparse footprint accounting,
the fully coupled affine recurrence, and the shifted Fibonacci specialization.

Finally, the accompanying Cubical Agda development checks the fixed model
theorem, the structural open-box theorem, a concrete SMod surface-adequacy
instance, recurrence arithmetic, lower-bound witnesses, and trace-accounting
discipline.  The development is evidence for the fixed theorem and for the
adequacy program, but the broad transfer claim remains conditional.  In
particular, this paper does not provide a parser or elaborator theorem for
arbitrary Cubical Agda programs, nor a theorem for every cubical foundation.

\subsection{The broader conjecture}

The fixed-calculus theorem is meant as evidence for a broader cubical principle.
Very roughly, cubical foundations with substitution-stable Kan composition and
filling, univalence, opaque sealing, and a faithful public-trace normalization
should exhibit the same pattern: unary trace is insufficient, binary trace is
necessary, and higher sealing-generated structural obligations are
horn-computationally derived.

This broader statement is intentionally left as a transfer program rather than
assumed as a theorem.  Different foundations package cubical composition,
univalence, strict equality, modalities, and module boundaries in different ways.
Each candidate system must supply its own account of what counts as a sealed
extension, which public fields are primitive, which fields are derived, and how
higher structural obligations are represented as open-box extension problems.
The abstract horn-computational interface in the final part of the paper records
exactly the hypotheses needed for such a transfer.

Thus the paper proves one instance carefully and proposes a route to others.  The
fixed theorem is the mathematical core; the adequacy interface is the boundary
between the proved result and the larger conjectural landscape.

\subsection{Roadmap}

Section~2 defines the sealed-extension calculus, the payload/trace distinction,
primitive and derived public fields, and the trace-cost invariant \(\mu\).
Section~3 develops normalized public signatures and the raw adequacy bridge,
including the SMod surface instance used as a running adequacy test.  Section~4
sets up structural horn obligations and the cubical operations used to compute
them.  Section~5 proves the horn-to-open-box theorem, open-box contractibility,
computational replacement, exact depth-two stabilization, recent-history
factorization, and the boolean-swap lower bound.  Section~6 derives the sparse
and full-coupling counting theorems, including the shifted Fibonacci endpoint.
Section~7 states the abstract horn-computational transfer interface and describes
the mechanized development and trusted boundary.  The appendices collect
contrasts, examples, and auxiliary formal details.
